\documentclass[11pt]{article}
\usepackage[margin=1in]{geometry}
\usepackage{booktabs}
\usepackage{hyperref}
\usepackage{xcolor}
\usepackage{listings}
\usepackage{tabularx}

\hypersetup{colorlinks=true, linkcolor=blue!60!black, urlcolor=blue!60!black}
\lstset{basicstyle=\ttfamily\small, breaklines=true, frame=single,
        backgroundcolor=\color{gray!8}}

\title{Interference-Aware Routing Schemes}
\author{Autonomous Networks Research Group (ANRG)\\University of Southern California}
\date{}

\begin{document}
\maketitle

\section{Overview}

When tasks in a DAG are placed on different nodes, data transfers traverse
multi-hop wireless paths where concurrent transmissions interfere with each
other.  The \textbf{routing scheme} decides \emph{which path} each transfer
takes and---in future work---\emph{when} it begins.  All schemes listed here
operate on top of HEFT task placement and the \texttt{csma\_bianchi} 802.11
interference model.

The schemes fall into three families:

\begin{table}[h]
\centering
\begin{tabular}{llp{5.5cm}}
\toprule
\textbf{Family} & \textbf{Schemes} & \textbf{When routes are decided} \\
\midrule
Baseline        & W, S         & Once, at topology load time (cached) \\
Static greedy   & GS, GC, GB, GO & Once, at DAG injection (after HEFT) \\
Dynamic greedy  & GSD            & At each \texttt{transfer\_start} event \\
Deferral \emph{(planned)} & GSD-D, GS-D, \ldots & At \texttt{transfer\_start}, may postpone \\
\bottomrule
\end{tabular}
\end{table}

%----------------------------------------------------------------------
\section{Baseline Schemes}

\subsection{W --- Widest-Path}

\begin{tabular}{@{}ll@{}}
\toprule
\textbf{Property} & \textbf{Value} \\
\midrule
Class        & \texttt{WidestPathRouting} \\
Objective    & Maximize bottleneck (min-link) bandwidth \\
Algorithm    & Modified Dijkstra with max-min semantics \\
Interference awareness & None---ignores concurrent traffic \\
\bottomrule
\end{tabular}

\medskip
Finds the path whose weakest link has the highest bandwidth.  Good when link
capacities vary widely (e.g., short vs.\ long hops have different PHY rates)
but oblivious to runtime contention.

\subsection{S --- Shortest-Path}

\begin{tabular}{@{}ll@{}}
\toprule
\textbf{Property} & \textbf{Value} \\
\midrule
Class        & \texttt{ShortestPathRouting} \\
Objective    & Minimize total path latency (sum of link latencies) \\
Algorithm    & Standard Dijkstra on link latencies \\
Interference awareness & None \\
\bottomrule
\end{tabular}

\medskip
Minimizes propagation delay.  In grids where all links have equal latency
this degenerates to minimum hop count.  Tends to concentrate traffic on short
paths, which can increase contention.

%----------------------------------------------------------------------
\section{Static Greedy Schemes (GS, GC, GB, GO)}

All four variants share the same two-stage architecture:

\begin{enumerate}\item \textbf{HEFT} decides task-to-node placement and produces timing estimates.
\item \textbf{Greedy route planner}
      (\texttt{InterferenceAwareRouting.plan\_routes()}) iterates over
      inter-node flows in a variant-specific order and, for each flow, picks
      the candidate path that maximizes \emph{total system throughput} across
      all concurrently-active flows (including previously routed ones).
\end{enumerate}

Routes are computed once at DAG injection and are fixed for the entire
execution.  The variants differ only in the order flows are considered.

\subsection{GS --- Greedy by Start Time}

\begin{tabular}{@{}lp{9cm}@{}}
\toprule
Flow ordering & Ascending estimated transfer start time \\
Rationale     & Route earlier flows first; later flows adapt around them \\
CLI           & \texttt{-{}-routing interference\_aware -{}-greedy-order start} \\
\bottomrule
\end{tabular}

\medskip
Earliest-first is the default and most natural ordering: it mirrors the
temporal sequence in which flows will actually appear during execution.

\subsection{GC --- Greedy by Criticality}

\begin{tabular}{@{}lp{9cm}@{}}
\toprule
Flow ordering & Descending HEFT upward-rank of the \emph{destination} task \\
Rationale     & Critical-path flows get first pick of low-interference routes \\
CLI           & \texttt{-{}-routing interference\_aware -{}-greedy-order criticality} \\
\bottomrule
\end{tabular}

\medskip
Upward rank (\texttt{ranku}) measures how much downstream work depends on a
task.  Routing the most critical transfers first gives them the best paths,
at the expense of less critical flows.

\subsection{GB --- Greedy by Bytes}

\begin{tabular}{@{}lp{9cm}@{}}
\toprule
Flow ordering & Descending data size \\
Rationale     & Large transfers dominate makespan; give them the best routes \\
CLI           & \texttt{-{}-routing interference\_aware -{}-greedy-order bytes} \\
\bottomrule
\end{tabular}

\medskip
When transfer sizes vary, the largest flows contribute most to makespan.
GB prioritizes them.  Less effective when all transfers are the same size
(as in many synthetic benchmarks).

\subsection{GO --- Greedy by Overlap}

\begin{tabular}{@{}lp{9cm}@{}}
\toprule
Flow ordering & Descending overlap degree (total temporal overlap with other flows) \\
Rationale     & Flows with the most contention should be routed first \\
CLI           & \texttt{-{}-routing interference\_aware -{}-greedy-order overlap} \\
\bottomrule
\end{tabular}

\medskip
Overlap degree is the sum of time-window overlaps between a flow and all
other flows (using HEFT timing estimates).  Flows competing for airtime with
many others get priority.

\subsection{Scoring Function (All Four Variants)}

For each candidate path, the planner computes:
\[
  \text{score} = \sum_{\text{concurrent flows}} \text{effective\_bottleneck\_bandwidth}
\]
where effective bandwidth accounts for:
\begin{itemize}\item \textbf{Interference factor} from \texttt{csma\_bianchi} (SINR
      degradation from hidden terminals + Bianchi MAC contention from
      conflict-graph neighbors)
\item \textbf{Per-link fair sharing} ($N$ flows on a link each get $1/N$ of
      the base bandwidth)
\end{itemize}
The candidate with the highest total-throughput score wins.

%----------------------------------------------------------------------
\section{Dynamic Greedy Scheme}

\subsection{GSD --- Greedy by Start Time, Dynamic}

\begin{tabular}{@{}lp{9cm}@{}}
\toprule
Class          & \texttt{DynamicInterferenceAwareRouting} \\
Objective      & Maximize \emph{this flow's} bottleneck bandwidth given actual state \\
When computed  & At each \texttt{transfer\_start} event \\
State used     & Actual set of active links and per-link transfer counts \\
CLI            & \texttt{-{}-routing interference\_aware\_dynamic} \\
\bottomrule
\end{tabular}

\medskip
Unlike the static variants (which estimate concurrency from HEFT timing),
GSD sees the \emph{real} network state at the moment a transfer begins.
This lets it:

\begin{itemize}\item React to transfers that finished early or late relative to HEFT estimates
\item See sibling transfers that started moments earlier at the same sim-time
\item Avoid paths that are congested \emph{right now}, not just predicted to be
\end{itemize}

\textbf{Trade-off:} GSD optimizes for \emph{this individual flow's}
bottleneck bandwidth rather than total system throughput.  It cannot
coordinate across flows the way the static planner can.

\textbf{Delegation:} During HEFT scheduling (before execution), GSD
delegates to \texttt{WidestPathRouting} for bandwidth estimates, since no
runtime state is available yet.

%----------------------------------------------------------------------
\section{Deferral (Planned---Not Yet Implemented)}

\subsection{Concept}

All current schemes start a transfer as soon as its source task completes.
\textbf{Deferral} adds an option: if starting a transfer \emph{now} would
severely degrade an already-active flow (or the new flow itself) due to
interference, the engine can \textbf{postpone} the new transfer until
conditions improve.

\subsection{Potential Variants}

\begin{description}\item[GSD-D (Dynamic + Deferral):] At \texttt{transfer\_start}, if the best
      available path scores below a threshold (e.g., ${<}\,50\%$ of the
      no-contention bandwidth), defer the transfer and re-evaluate after the
      next \texttt{transfer\_complete} event.
\item[GS-D (Static + Deferral):] Use GS route planning but allow the
      execution engine to hold back flows whose planned route is currently
      congested.
\end{description}

\subsection{Key Design Questions}

\begin{enumerate}\item \textbf{Deferral trigger:} Score threshold?  Comparison to
      no-contention baseline?  Always defer if any link on the path is above
      $N$ concurrent flows?
\item \textbf{Re-evaluation:} On every \texttt{transfer\_complete}?  After a
      fixed timeout?
\item \textbf{Starvation:} How to guarantee deferred transfers eventually
      execute (maximum deferral budget, priority escalation)?
\item \textbf{Impact on HEFT:} Deferral changes actual timing relative to
      HEFT estimates.  Does this cascade and hurt downstream tasks?
\end{enumerate}

\subsection{Expected Benefit}

Deferral is most useful when:
\begin{itemize}\item A few high-bandwidth transfers dominate makespan and would complete
      faster if not forced to share with small concurrent transfers
\item The DAG has enough parallelism that deferring one transfer doesn't
      immediately stall a critical-path task
\end{itemize}

%----------------------------------------------------------------------
\section{Configuration Reference}

\subsection{CLI}

\begin{lstlisting}
--routing {direct,widest_path,shortest_path,
           interference_aware,interference_aware_dynamic}
--greedy-order {start,criticality,bytes,overlap}
--routing-hop-cutoff N
--interference {none,proximity,csma_clique,csma_bianchi}
\end{lstlisting}

\subsection{YAML}

\begin{lstlisting}
config:
  routing: interference_aware
  greedy_order: start
  routing_hop_cutoff: 4
  interference: csma_bianchi
\end{lstlisting}

\subsection{Shared Parameters}

\begin{tabular}{@{}lcp{7cm}@{}}
\toprule
\textbf{Parameter} & \textbf{Default} & \textbf{Effect} \\
\midrule
\texttt{hop\_cutoff}     & 4  & Maximum hops in candidate path enumeration \\
\texttt{max\_candidates} & 20 & Maximum candidate paths evaluated per flow \\
\bottomrule
\end{tabular}

%----------------------------------------------------------------------
\section{Summary}

\begin{table}[h]
\centering
\small
\begin{tabular}{lllll}
\toprule
\textbf{Scheme} & \textbf{Route timing} & \textbf{Objective} & \textbf{Flow ordering} & \textbf{Interference input} \\
\midrule
W    & Load time       & Max bottleneck BW       & ---              & None \\
S    & Load time       & Min total latency        & ---              & None \\
GS   & DAG inject      & Max system throughput     & Est.\ start time & HEFT estimates \\
GC   & DAG inject      & Max system throughput     & HEFT upward rank & HEFT estimates \\
GB   & DAG inject      & Max system throughput     & Data size (desc) & HEFT estimates \\
GO   & DAG inject      & Max system throughput     & Overlap degree   & HEFT estimates \\
GSD  & Transfer start  & Max flow bottleneck BW    & ---              & Actual active links \\
GSD-D & Transfer start & Max flow bottleneck BW    & ---              & Actual + deferral \\
\bottomrule
\end{tabular}
\end{table}

\end{document}
